\begin{ledgroupsized}[r]{120mm}
\footnotesize
\pstart
\noindent\textbf{Überlieferung:}
\pend
\end{ledgroupsized}
\begin{ledgroupsized}[r]{114mm}
\footnotesize
\pstart \parindent -6mm
\makebox[6mm][l]{\textit{L}}%
Konzept:
LH~XXXV~11,~14 Bl.~1\textendash2.
Ein Bogen~4\textsuperscript{o}.
\href{\Scan{LH_35_11_14_001-002wz}}%
{Wasserzeichen}
im Falz.
Vier Seiten.
\pend
\end{ledgroupsized}
\begin{ledgroupsized}[r]{114mm}
\footnotesize
\pstart
\parindent -6mm
\makebox[6mm][l]{\textit{E\textsuperscript{1}}}%
Erstdruck: \cite{01023}\textit{AE} November 1685, S.\;501\textendash505;
Abbildungen (\textit{Tab.~XIII, Fig. 2\textendash3}) ebd. zw. S.\;494 u. 495.
(Unsere Druckvorlage.)
Nachdr. \textsc{Lamarra/Palaia} 2005, S.\;84\textendash89.\cite{01289}
Es existieren verschiedene, in der Interpunktion und im Seitenumbruch von S.\;504 auf S.\;505 abweichende Druckfassungen der \textit{Acta Eruditorum}.
\pend%
\end{ledgroupsized}
\begin{ledgroupsized}[r]{114mm}
\footnotesize 
\pstart \parindent -6mm
\makebox[6mm][l]{\textit{E\textsuperscript{2}}}%
Abdruck (nach \textit{E\textsuperscript{1}}): 
\cite{01421}\textit{Opuscula omnia Actis eruditorum Lipsiensibus inserta, quae ad universam mathesin, physicam, medicinam, anatomiam, chirurgiam et philologiam pertinent}, Venedig 1740--1743,  Bd.~I, S.\;358--362; Abbildungen am Ende des Bandes.
\pend
\end{ledgroupsized}
\begin{ledgroupsized}[r]{114mm}
\footnotesize 
\pstart \parindent -6mm
\makebox[6mm][l]{\textit{E\textsuperscript{3}}}%
Abdruck (nach \textit{E\textsuperscript{1}}): 
G.\,W. \textsc{Leibniz}, \textit{Opera omnia}, hrsg. von L.~\textsc{Dutens}, 
Genf 1768, Bd.~III, %
S.\;175\textendash179; 
Abbildungen am Ende des Bandes, Tab.~VII, Fig. 32\textendash33.\cite{00150}
\pend
\end{ledgroupsized}
\begin{ledgroupsized}[r]{114mm}
\footnotesize 
\pstart \parindent -6mm
\makebox[6mm][l]{\textit{E\textsuperscript{4}}}%
Abdruck (nach \textit{E\textsuperscript{1}}): 
G.\,W. \textsc{Leibniz}, \textit{Mathematische Schriften}, hrsg. von C.\,I. \textsc{Gerhardt}, 
Bd.~VI, Halle 1860, S.\;112\textendash117; 
Abbildungen auf dem ersten Faltblatt am Ende des Bandes, Fig.~9\textendash10.\cite{01043}
\pend
\end{ledgroupsized}
\begin{ledgroupsized}[r]{114mm}
\footnotesize 
\pstart
\parindent -6mm
\makebox[6mm][l]{\textit{Ü}}%
(ins Französische): 
G.\,W. \textsc{Leibniz}, \textit{Oeuvre concernant la physique}, hrsg. von P.~\textsc{Peyroux}, 
Paris 1985, S.\;21\textendash24; 
Abbildungen im Text.\cite{01248}
\pend
\end{ledgroupsized}
\vspace{8mm}
\selectlanguage{latin}
\pstart
\noindent
\lbrack%
\textit{Nachfolgender Text bis \refpassage{AE_1685_501-505_16a}{AE_1685_501-505_16a} 
nur in \textit{L}, fehlt in E\textsuperscript{1}:}\rbrack
\pend
\vspace{0.5em}
\pstart%
\noindent
\href{\Scan{LH_35_11_14_001r}}%
{\lbrack1~r\textsuperscript{o}\rbrack} %
Ad
\edtext{brevem annotationem in \cite{01023}\textit{Actis} Novemb. 1684 p.~514}{%
\lemma{brevem annotationem \lbrack...\rbrack\ p.~514}%
\Cfootnote{%
\indexN{\textls{Vanni}, Giovanni Francesco 1638\textendash1709}Vannis anonymes 
\cite{01995}\textit{Specimen libri de momentis gravium}
(\cite{01023}\textit{AE} November 1684, S.\;511\textendash513)
war von einer ebenfalls anonymen, titellosen Bemerkung Leibnizens begleitet
(ebd. S.\;514),
die unter dem Titel \glqq Brevis annotatio circa Specimen libri de momentis gravium\grqq\ 
in diesem Band als N.\;\ref{55606} erscheint.}}
\indexN{\textls{Vanni}, Giovanni Francesco 1638\textendash1709}autor 
\edtext{\cite{01995}\textit{Speciminis de momentis gravium}}{%
\lemma{\textit{Speciminis de momentis gravium}}%
\Cfootnote{\indexN{\textls{Vanni}, Giovanni Francesco 1638\textendash1709}\textsc{G.\,F. Vanni},
\cite{01995}\textit{Specimen}.}} 
\textls{respondit},
\edtext{ut intelligo:}{%
\lemma{ut intelligo}%
\Cfootnote{%
Leibniz hatte über 
\indexN{\textls{Magliabechi} (Magliabecchi, Magliabecki, Magliabekius, Magliabequius), Antonio 1633\textendash1714}A.~Magliabechi und 
\indexN{\textls{Pfautz} (Pfauzius), Christoph 1645\textendash1711}C.~Pfautz
einen
Brief \indexN{\textls{Vanni}, Giovanni Francesco 1638\textendash1709}Vannis erhalten (siehe \cite{02009}\textit{LSB} III, 4 N.\;96) und daraus exzerpiert (ebd. \cite{02010}N.\;97).
Beide nachfolgenden Einwände sind diesem Brief entnommen; Leibniz hatte sie bereits in den Konzepten N.\;\ref{32310} und N.\;\ref{32313} (dort ab \refpassage{LBr_724_018-019_8b}{LBr_724_018-019_8b}), zwei Vorarbeiten zum vorliegenden Aufsatz, wiedergegeben und beantwortet.}}
\textls{primo} instantiae objectae hic
non esse locum contra se, quia in ea anguli planorum
non sint recti,\indexS{angulus planorum non rectus} in
\edtext{qu\lbrack o\rbrack}{%
\lemma{}%
\Bfootnote{qua \textit{L ändert Hrsg.}}}
consistat nervus suae demonstrationis\lbrack;\rbrack\
\textls{secundo} frustra
\edtext{negari momenta}{%
\lemma{negari}%
\Bfootnote{%
\textit{(1)}~plana in %
\textit{(2)}~momenta %
\textit{L}}} 
in planis
declivibus\indexS{planum declive} facientibus angulum rectum\indexS{plana declivia angulum rectum facientia} 
in unum posse addi,\indexS{additio momentorum} ut componant momentum totale.\indexS{momentum totale} 
Nam evidens esse, in unum posse addi momenta\indexS{additio momentorum} non secus ac pondera,
\indexS{additio ponderum} quibus illa aequari demonstratum sit.
\pend
\pstart
\textls{Replicari} posse videtur ad \textls{primum},
equidem cessare instantiam, aut saltem esse immutandam,
si \indexN{\textls{Vanni}, Giovanni Francesco 1638\textendash1709}autor 
\cite{01995}\textit{Speciminis} additionem momentorum\indexS{additio momentorum} 
ad solum angulum rectum planorum\indexS{angulus planorum rectus} restringat.
\pend
\pstart
Ad \textls{secundum}
\edtext{(quod ipsam praetensae ab autore demonstrationis solutionem concernit)}{%
\lemma{}%
\Bfootnote{%
(quod %
\textit{(1)}~ab %
\textit{(2)}~ipsam praetensae %
\textit{(a)}~\bars{ab aut \textit{erg.~u.~gestr.}} demonstrationis solutionem %
\textit{(b)}~ab autore %
\textit{(aa)}~solutionis demonstrationem concernit)
\textit{(bb)}~demonstrationis solutionem concernit)
\textit{erg.~L}}}
replicatur, negatum
non esse momenta\indexS{additio momentorum quaecunque} quaecunque in unum posse addi, sed ab autore in
unum addita aequari momento totali.\indexS{momentum totale} Negatum quoque non est
momenta quibus singula plana revera premuntur,\indexS{momentum quo planum singulum premitur} in unum
addita\indexS{additio momentorum} componere totale seu absolutum momentum globi,\indexS{momentum totale seu absolutum} nam
id quoque per se patet, cum tota vis globi\indexS{vis globi} a planis sustineatur\lbrack;\rbrack\
sed hoc tantum negatum est, ea momenta\indexS{momentum gravis planum propellendi} quibus grave planum
alterutrum\lbrack,\rbrack\ si cedere posset, propelleret\lbrack,\rbrack\ hoc est (in casu recti
planorum anguli)\indexS{angulus planorum rectus} momenta quibus in altero plano
\edtext{declivi si posset descenderet,\indexS{momentum gravis in plano descendendi}}{%
\lemma{declivi}%
\Bfootnote{%
\textit{(1)}~descenderet, si posset %
\textit{(2)}~si posset descenderet, %
\textit{L}}}
coincidere cum his quibus ipsa plana in casu impediti
descensus\indexS{descensus impeditus} premuntur, ac proinde in unum addita\indexS{additio momentorum}
\edtext{\lbrack momentum\rbrack}{%
\lemma{}%
\Bfootnote{planum \textit{L ändert Hrsg.}}}
totale\indexS{momentum totale} conficere.
\pend 
\pstart
Ut autem utilitatem aliquam ex hac concertatione capiat lector,
nec tantum in refutando, aut refutationi
\edtext{communis ac ut putamus verae sententiae}{%
\lemma{communis}%
\Bfootnote{%
\bars{%
\textit{(1)}~ac ve %
\textit{(2)}~ac ut putamus verae
\textit{erg.}}
sententiae~\textit{L}}}
respondendo, occupemur\lbrack,\rbrack\ placet primum demonstrationem
\edtext{aliquam nostram brevem et evidentem nec ut putamus inelegantem Regulae}{%
\lemma{aliquam}%
\Bfootnote{%
\bars{nostram \textit{erg.}} %
brevem %
\textit{(1)}~, ut %
\textit{(2)}~et evidentem nec ut putamus inelegantem %
\textit{(a)}~communis %
\textit{(b)}~Regulae %
\textit{L}}}
Staticorum\indexS{Statici} circa momenta inclinata\indexS{momentum inclinatum} receptae\indexS{regula Staticorum recepta} subjicere, %
\edtext{neque enim ad manus sunt demonstrationes 
\indexN{\textls{Galilei} (Galilaeus, Galileus), Galileo 1564\textendash1642}Galilaei et 
\indexN{\textls{Torricelli} (Torricellius), Evangelista 1608\textendash1647}Torricellii, quas doctissimus 
\indexN{\textls{Vanni}, Giovanni Francesco 1638\textendash1709}autor 
\cite{01995}\textit{Speciminis} sibi satisfacere negat.%
}{%
\lemma{neque enim \lbrack...\rbrack\ negat}%
\Cfootnote{%
\textsc{Vanni},\indexN{\textls{Vanni}, Giovanni Francesco 1638\textendash1709} \cite{01995}\textit{Specimen}, hier S.\;511.
Siehe
\indexN{\textls{Galilei} (Galilaeus, Galileus), Galileo 1564\textendash1642}%
\textsc{G.~Galilei}, %
\cite{00050}\textit{Discorsi}, %
Giornata Terza, De motu naturaliter accelerato; \glqq Aggiunta postuma dell'Autore\grqq\ %
nach dem Scholion zum 2.~Theorem (der Zusatz erschien erst in der posthumen Werkausgabe: \cite{01084}\textit{Opere}, Bologna 1656, Bd.~II, S.\;132\lbrack ter\rbrack\textendash134; \cite{00048}\textit{GO} VIII, S.\;216\textendash219 im Apparat) %
und 
\indexN{\textls{Torricelli} (Torricellius), Evangelista 1608\textendash1647}%
\textsc{E.~Torricelli}, \textit{De motu gravium}, lib.~I
(\cite{00304}\textit{Opera geometrica}, Florenz 1644, S.\;99;
\cite{00108}\textit{TO}~II, S.\;105).}}
Huic demonstrationi 
\edtext{deinde \href{\Scan{LH_35_11_14_001v}}%
{\lbrack1~v\textsuperscript{o}\rbrack} subnectemus tentatam Solutionem}{%
\lemma{deinde}%
\Bfootnote{%
\textit{(1)}~solvemus %
\textit{(2)}~\href{\Scan{LH_35_11_14_001v}}%
{\lbrack1~v\textsuperscript{o}\rbrack} subjiciemus %
\textit{(3)}~subnectemus
\textit{(a)}~modum quo putamus unumquodque planum %
\textit{(b)}~tentatam Solutionem %
\textit{L}}}
casus ab \indexN{\textls{Vanni}, Giovanni Francesco 1638\textendash1709}autore \cite{01995}\textit{Speciminis} propositi,
\edtext{seu determinationem}{%
\lemma{seu}%
\Bfootnote{%
\textit{(1)}~modum quo ex %
\textit{(2)}~momentum quo utrumque planorum inclinatorum, quant %
\textit{(3)}~determinationem %
\textit{L}}}
momenti\indexS{momentum quo planum premitur} quo unumquodque planorum propositorum
angulum rectum facientium\indexS{plana angulum rectum facientia}
\edtext{premitur\lbrack,\rbrack\ \lbrack Q\rbrack uoniam}{%
\lemma{}%
\Bfootnote{premitur. Quoniam \textit{L ändert Hrsg.}}} 
id nondum videmus explicatum\lbrack;\rbrack\
ut ita Clarissimum \indexN{\textls{Vanni}, Giovanni Francesco 1638\textendash1709}autorem 
\edlabel{AE_1685_501-505_8a}%
\edtext{}{%
\xxref{AE_1685_501-505_8a}{AE_1685_501-505_8b}%
\lemma{\textit{Speciminis}}%
\Bfootnote{%
\textit{(1)}~aliosque, %
\textit{(2)}~pariter ac Adm.\,R.\,P.~Kochanskium cujus in \textit{Actis} Junii hujus anni elegantes meditationes habentur,~\textit{L}}}%
\cite{01995}\textit{Speciminis} pariter ac %
\edtext{Adm.\,R.\,P.~Kochanskium\indexN{\textls{Kochański}, Adam Adamandus 1631\textendash1700} cujus in \textit{Actis} Junii hujus anni elegantes meditationes habentur,\edlabel{AE_1685_501-505_8b}}{%
\lemma{Adm. \lbrack...\rbrack\ habentur}%
\Cfootnote{%
\textsc{A.\,A. Kochański}, %
\cite{01996}\textit{Consideratio speciminis libri de momentis gravium, Authore J.\,F.\,V. Lucensi}, %
\cite{01023}\textit{AE} Juni 1685, %
S.\;262\textendash265.}}
ad argumentum hoc porro illustrandum
\edlabel{AE_1685_501-505_1a}%
\edtext{}{%
\xxref{AE_1685_501-505_1a}{AE_1685_501-505_1b}%
\lemma{}%
\Bfootnote{%
invitemus.
\textls{Demonstratio Geometrica regulae apud Staticos
receptae de momentis gravium in planis inclinatis} %
\bars{et solutio casus %
\textit{(1)}~de gra %
\textit{(2)}~in \textit{Actis} Novembr.~1684 p.~512 propositi de %
\textit{(a)}~gravi %
\textit{(b)}~globo %
duobus planis angulum rectum facientibus simul incumbente, quantum %
unumquodque planorum premat autore G.\,G.\,L. \textit{erg.}}%
~\textit{L}
\textbullet~\lbrack S.\;501\rbrack\
G.\,G.\,L. \lbrack...\rbrack\ determinans.~\textit{E\textsuperscript{1}}}}%
invitemus.\edlabel{AE_1685_501-505_16a}	%
\pend
\normalsize 
\vspace{0.5\baselineskip}
\pstart
\centering
\lbrack S.\;501\rbrack\  %
\hfill
\edlabel{AE_1685_501-505_16b}%
G.\,G.\,L. DEMONSTRATIO GEOMETRICA
\hfill
\phantom{\lbrack S.\;501\rbrack}
\linebreak
regulae apud Staticos\indexS{Statici} receptae\indexS{regula apud Staticos recepta}
\textls{de momentis gravium in planis inclinatis,}\indexS{momentum gravis in plano inclinato}
nuper in dubium vocatae: et solutio casus elegantis,
\edtext{in \textit{Actis} Novembr.~1684 pag.~512}{%
\lemma{in \textit{Actis} Novembr.~1684 pag.~512}%
\Cfootnote{\textsc{Vanni},\indexN{\textls{Vanni}, Giovanni Francesco 1638\textendash1709} \cite{01995}\textit{Specimen}, bes. S.\;512.}}
propositi, \textls{de globo\indexS{globus duobus planis simul incumbens} duobus planis angulum rectum facientibus\indexS{plana angulum rectum facientia} simul incumbente;}\linebreak%
quantum unumquodque planorum prematur, determinans.\indexS{momentum quo planum inclinatum premitur}\edlabel{AE_1685_501-505_1b}
\pend 
\begin{edfigure}%
\hfill
\hspace*{2mm}
\includegraphics[scale=1]{%
gesamttex/edit_VIII-4/images/LH_35_11_14_001-002_d_001v.pdf}
\hspace*{15mm}%
\includegraphics[scale=1]{%
gesamttex/edit_VIII-4/images/AE_1685_501-505_d1.pdf}
\hfill
\edcaption{%
\edtext{%
\lbrack\textit{Fig.~1a; L, Bl.~1~v\textsuperscript{o}}\rbrack}{\lemma{\lbrack\textit{Fig.~1a}\rbrack}\Cfootnote[nonum]{%
Tlw. Blindzeichnung.}}
\hspace*{32.5mm}
\edtext{\lbrack\textit{Fig.~1b; E\textsuperscript{1}}\rbrack}{\lemma{\lbrack\textit{Fig.~1b}\rbrack}\Cfootnote[nonum]{%
Die Verbindungsstrecke zwischen \textit{D} und \textit{E} ändert Hrsg. nach ihrer Definition im Text.}}%
\hspace*{7.5mm}}
\end{edfigure}
\pstart
\noindent
\edtext{Sit in 
\edtext{\textls{Fig. 2}}{%
\lemma{\textls{Fig. 2}}%
\Cfootnote{Siehe hier das Diagramm \lbrack\textit{Fig.~1b}\rbrack.}}
planum verticale\indexS{planum verticale} \textit{ADB}, et inclinatum\indexS{planum inclinatum} \textit{ACE}; et}{%
\lemma{Sit}%
\Bfootnote{%
\textit{(1)}~planum verticale \bars{in figur \textit{erg.}} \textit{ADB}, %
\textit{(2)}~\bars{in fig. 1 \textit{erg.}} planum verticale \textit{ADB}, %
\textit{(a)}~\bars{et \textit{erg.}} inclinatum %
\textit{(b)}~et inclinatum \textit{ACE} et%
~\textit{L}
\repeatlemma \lbrack...\rbrack\ \textit{ACE}; et~\textit{E\textsuperscript{1}}}}
in illo descendere tendat grave\indexS{grave descendere tendens} \textit{B}, in hoc
\edtext{grave \textit{C}, quae designo per puncta in quae incidunt eorum centra gravitatis.\indexS{centrum gravitatis} Sint autem}{%
\lemma{grave \textit{C},}%
\Bfootnote{%
sint autem%
~\textit{L}
\repeatlemma \lbrack...\rbrack\ autem ~\textit{E\textsuperscript{1}}%
}}
gravia connexa\indexS{gravia fune connexa}
\edtext{fune\indexS{funis} \textit{BAC}, et in aequilibrio, angulusque \textit{ABC} rectus. \textls{Ajo fore grave \textit{B} ad grave \textit{C}, ut recta \textit{AB} ad rectam \textit{AC}}, quod sic}{%
\lemma{fune \textit{BAC}}%
\Bfootnote{%
\bars{et in aequilibrio \textit{erg.}},
angulusque \textit{ABC} sit rectus
\textit{(1)}~, et gravia eam inter se rationem habeant, ut sint in %
\textit{(a)}~aequilibrio, %
\textit{(b)}~aequilibrio %
\textit{(2)}~. Ajo fore grave \textit{B} ad grave \textit{C}, ut recta \textit{AB} ad rectam \textit{AC}. Quod sic%
~\textit{L}
\repeatlemma \lbrack...\rbrack\ sic~\textit{E\textsuperscript{1}}%
}}
ostendo: 
\edlabel{AE_1685_501-505_15a}%
Funem\indexS{funis gravitate carens}, (qui gravitate carere
\edtext{intelligitur),}{%
\lemma{intelligitur)}%
\Bfootnote{\textit{L}
\repeatlemma,~\textit{E\textsuperscript{1}}}}
manu prehendendo
\edtext{in \textit{C},}{%
\lemma{in \textit{C}}%
\Bfootnote{\textit{L}
\repeatlemma~,~\textit{E\textsuperscript{1}}}}
trahendoque grave \textit{C} deorsum in \textit{E}, adeoque simul grave \textit{B} sursum in \textit{D} 
(quo facto aequales erunt \textit{CE} et \textit{BD}) 
\edtext{aequilibrium\indexS{aequilibrium},}{%
\lemma{aequilibrium}%
\Bfootnote{\textit{L}
\repeatlemma,~\textit{E\textsuperscript{1}}}}
quod prius erat inter gravia, 
\edtext{manere, nec manum ab ipsis juvari vel impediri, nec}{%
\lemma{manere,}%
\Bfootnote{%
\bars{nec manum ab illis 
\textit{(1)}~juv %
\textit{(2)}~juvari vel impediri, %
\textit{erg.}} nec~\textit{L}
\repeatlemma \lbrack...\rbrack\ nec~\textit{E\textsuperscript{1}}}}
proinde eorum centrum gravitatis\indexS{centrum gravitatis commune} commune 
\edtext{(seu centrum totius ex ipsis compositi, sive centrum aequilibrii\indexS{centrum aequilibrii})}{%
\lemma{}%
\Bfootnote{%
(seu \lbrack...\rbrack\ aequilibrii) %
\textit{fehlt~L}}}
attolli vel
\edtext{deprimi.\edlabel{AE_1685_501-505_15b}	%
Et cum}{%
\lemma{deprimi}%
\Bfootnote{%
\textit{(1)}~. Cumque %
\textit{(2)}~: itaque cum %
gravium \bars{in \textit{streicht Hrsg.}} %
situ \textit{B}, \textit{C}, centrum gravitatis esset in recta \textit{BC}, horizontali, %
in ea manebit, non obstante hac gravium translatione quae ipsa in aequilibrio relinquit %
\textit{(3)}~.~Cumque gr
\textit{(4)}~.~\bars{Et~\textit{erg.}} cum~\textit{L}}}
gravium in situ \textit{B}, \textit{C}, centrum gravitatis\indexS{centrum gravitatis commune} commune esset in recta
\edtext{\textit{BC}, horizonti parallela,\indexS{recta horizonti parallela} ideo}{%
\lemma{\textit{BC},}%
\Bfootnote{%
\bars{horizonti parallela \textit{erg.}} %
\textit{(1)}~eoru %
\textit{(2)}~ideo %
\textit{L}
\repeatlemma horizonti parallela, ideo~\textit{E\textsuperscript{1}}}}
eorundem centrum\indexS{centrum gravitatis}
\edtext{in situ \lbrack\textit{D}, \textit{E},\rbrack\ positorum}{%
\lemma{in situ}%
\Bfootnote{\textit{D}, \textit{E},~\textit{L}
\repeatlemma \bars{\textit{DE} \textit{ändert Hrsg. nach L}} positorum~\textit{E\textsuperscript{1}}}}
manebit in
\edtext{recta \textit{BC}, alioqui attolleretur vel deprimeretur. At idem est et in recta \textit{DE}; quippe quae eorum centra gravitatis\indexS{centrum gravitatis}, \textit{B} et \textit{C}, jungit: erit ergo in \textit{H} puncto communi}{%
\lemma{recta}%
\Bfootnote{%
\textit{BC}; at idem est \bars{et~\textit{erg.}} %
in recta \textit{DE}. Erit ergo in
\textit{(1)}~intersec %
\textit{(2)}~\textit{H} puncto %
\textit{(a)}~intersectionis %
\textit{(b)}~communi~\textit{L}
\repeatlemma \lbrack...\rbrack\ communi~\textit{E\textsuperscript{1}}}}
rectarum \textit{BC} et \textit{DE}. Jam ex Geometria\indexS{Geometria} ostendi potest 
\textls{Lemma} facile nec
\edtext{injucundum:}{%
\lemma{injucundum}%
\Bfootnote{,~\textit{L}
\repeatlemma:~\textit{E\textsuperscript{1}}}}
si sint \textit{CE} et \textit{BD} aequales, fore \textit{DH} ad
\edtext{\textit{HE},}{%
\lemma{\textit{HE}}%
\Bfootnote{\textit{L}
\repeatlemma,~\textit{E\textsuperscript{1}}}}
ut \textit{AC} ad
\edtext{\textit{AB},}{%
\lemma{\textit{AB}}%
\Bfootnote{.~\textit{L}
\repeatlemma,~\textit{E\textsuperscript{1}}}}
(\textls{de quo mox}.) Est autem etiam grave \textit{E} (sive \textit{C}) ad grave \textit{D} (sive \textit{B}) ut \textit{DH} ad \textit{HE}, quia \textit{H} est centrum
\edtext{aequilibrii\indexS{centrum aequilibrii} ipsorum. Ergo}{%
\lemma{aequilibrii}%
\Bfootnote{eorum. Ergo~\textit{L}
\repeatlemma ipsorum. Ergo~\textit{E\textsuperscript{1}}}}
erit grave \textit{C} ad
\edtext{grave \textit{B},}{%
\lemma{grave}%
\Bfootnote{\textit{B}~\textit{L}
\repeatlemma \textit{B}~,~\textit{E\textsuperscript{1}}}} 
ut \textit{AC} ad
\edtext{\textit{AB},}{%
\lemma{\textit{AB}}%
\Bfootnote{\textit{L}
\repeatlemma,~\textit{E\textsuperscript{1}}}}
quod asserebatur.
\pend
\pstart
\textls{Idem aliter etiam sine assumto lemmate} 
\edtext{illo geometrico ostendi potest, si}{%
\lemma{illo}%
\Bfootnote{Geometrico ostendi potest si~\textit{L}
\repeatlemma geometrico ostendi potest, si~\textit{E\textsuperscript{1}}}}
grave \textit{B} attollatur usque ad \textit{A}, seu si
\edtext{\textit{(D)} incidat in ipsum punctum \textit{A}, eo enim casu \textit{C}\textit{(E)} seu \textit{B}\textit{(D)}}{%
\lemma{\textit{(D)}}%
\Bfootnote{%
\textit{(1)}~incidet %
\textit{(2)}~incidat in ipsum punctum \textit{A}, %
\textit{(a)}~quo %
\textit{(b)}~eo enim casu \textit{C}(\textit{E}) %
\textit{(aa)}~aequalis %
\textit{(bb)}~seu \textit{B}(\textit{D})~\textit{L}}}
erit aequalis \textit{BA}, et \textit{(H)} incidet in \textit{C}, et \textit{(H)}\textit{(D)} co\lbrack S.\;502\rbrack incidet cum
\edtext{\textit{CA}; eritque \textit{C} centrum aequilibrii\indexS{centrum aequilibrii} gravium 
ex \textit{B} et \textit{C} translatorum in \textit{A} et \textit{(E)}, ac proinde grave in \textit{(E)} 
hoc est grave \textit{C}, ad grave in \textit{A}, hoc est grave \textit{B}, erit}{%
\lemma{\textit{CA}}%
\Bfootnote{, eritque \textit{C} gravium
\textit{(1)}~\textit{B} et \textit{C} %
\textit{(2)}~ex \textit{B} \lbrack...\rbrack\ et \textit{(E)}, centrum %
\textit{(a)}~gravitatis %
\textit{(b)}~aequilibrii, ac proinde %
\bars{grave in \textit{erg.}} \textit{(E)} %
\textit{(aa)}~sive \textit{C} ad \textit{B} sive \textit{A} %
\textit{(bb)}~hoc est grave \textit{C} ad grave in \textit{A} hoc est grave \textit{B} erit~\textit{L}
\repeatlemma; \lbrack...\rbrack\ erit~\textit{E\textsuperscript{1}}}}
ut \textit{AC} ad \textit{C}\textit{(E)} %
\edlabel{AE_1685_501-505_2a}%
\edtext{}{%
\xxref{AE_1685_501-505_2a}{AE_1685_501-505_2b}%
\lemma{sive}%
\Bfootnote{%
\textit{AB}~. Placet%
~\textit{L}
\repeatlemma \textit{AB}, ut ante. Placet~\textit{E\textsuperscript{1}}}}%
sive \textit{AB}, ut ante.
\pend
\pstart
Placet\edlabel{AE_1685_501-505_2b}
tamen ipsius propositionis Geometricae supradictae demonstrationem 
\edlabel{AE_1685_501-505_7a}%
\edtext{}{%
\xxref{AE_1685_501-505_7a}{AE_1685_501-505_7b}%
\lemma{subjicere.}%
\Bfootnote{%
\textls{Lemma: omnium triangulorum angulum coincidentem et}~\textit{L}
\repeatlemma \lbrack...\rbrack\ et~\textit{E\textsuperscript{1}}}}%
subjicere.
\pend
\pstart
LEMMA: \textls{omnium triangulorum, angulum coincidentem, et\edlabel{AE_1685_501-505_7b} 
summam laterum circa eundem angulum aequalem habentium, bases}
\edtext{\textls{a basi unius ex ipsis trianguli in eadem}}{%
\lemma{\textls{a basi}}%
\Bfootnote{%
\textit{(1)}~\textls{ejus trianguli} %
\textit{(2)}~\textls{unius ex ipsis trianguli} %
\textit{(a)}~, \textls{quod inter ipsa rectangulum est} %
\textit{(b)}~\textls{in eadem}~\textit{L}}}
\textls{ratione}
\edtext{\textls{secantur;}}{%
\lemma{\textls{secantur}}%
\Bfootnote{\,\textit{L}
\repeatlemma;~\textit{E\textsuperscript{1}}}}
\textls{quae scilicet est ratio}
\edtext{\textls{laterum hujus trianguli positorum}}{%
\lemma{\textls{laterum}}%
\Bfootnote{%
\textit{(1)}~\textls{trianguli illius} %
\textit{(2)}~\textls{hujus trianguli} %
\textit{(a)}~\textls{rectanguli} %
\textit{(b)}~\textls{positorum}~\textit{L}}}
\textls{circa angulum}
\edlabel{AE_1685_501-505_6a}%
\edtext{}{%
\xxref{AE_1685_501-505_6a}{AE_1685_501-505_6b}%
\lemma{\textls{communem.}}%
\Bfootnote{%
\href{\Scan{LH_35_11_14_002r}}%
{\lbrack2~r\textsuperscript{o}\rbrack} Sit in \textls{eadem} 
\textit{(1)}~\textls{figura} 
\textit{(2)}~\textls{figur.~1} angulus communis \textit{FAM}, et triangula quotcunque \textit{ABC},~\textit{L}
\repeatlemma \lbrack...\rbrack\ ut \textit{ABC},~\textit{E\textsuperscript{1}}}}%
\textls{communem.}
\pend
\pstart
Sit angulus communis \textit{FAM} \textls{in eadem} 
\edtext{\textls{Figur.~2,}}{%
\lemma{\textls{Figur.~2}}%
\Cfootnote{Siehe hier das Diagramm \lbrack\textit{Fig.~1b}\rbrack.}}
et triangula quotcunque ut \textit{ABC},\edlabel{AE_1685_501-505_6b}
\textit{ADE} angulum \textit{BAC}, \textit{DAE} habentia coincidentem, et summam laterum
$BA+AC$ vel $DA+AE$
semper aequalem (unde et \textit{BD} aequalis erit ipsi \textit{CE}) unumque ex his
\edtext{triangulis assumatur ut \textit{ABC}: dico a basi ejus \textit{BC} secari basin \textit{DE} trianguli alterius cujuscunque supradictas conditiones habentis \textit{ADE}, in partes}{%
\lemma{triangulis}%
\Bfootnote{%
\textit{(1)}~ut \textit{ABC} sit rectangulum %
\textit{(2)}~assumatur ut \textit{ABC}. Dico a basi ejus \textit{BC}
\textit{(a)}~, basin %
\textit{(aa)}~alterius trianguli, \textit{DE}, %
\textit{(bb)}~\textit{DE} %
\textit{(aaa)}~al %
\textit{(bbb)}~trianguli \bars{alterius cujuscunque supradictas conditiones habentis \textit{erg.}} 
\textit{ADE} secari in partes
\textit{(c)}~secari basin \lbrack...\rbrack\ \textit{ADE} in partes~\textit{L}
\repeatlemma \lbrack...\rbrack\ in partes~\textit{E\textsuperscript{1}}}}
\textit{DH},
\edtext{\textit{HE},}{%
\lemma{\textit{HE}}%
\Bfootnote{\textit{L}
\repeatlemma,~\textit{E\textsuperscript{1}}}}
quae sint in ratione \textit{AC}, \textit{AB}. Nimirum ex \textit{D} ad \textit{AC} ducatur \textit{DG} parallela ipsi \textit{BC}, et ex \textit{G} ad \textit{BC} ducatur
\edtext{\textit{GL}, parallela ipsi \textit{AB}. Jam (ob triangula \textit{ABC}, \textit{GLC}, similia) est \textit{AC} ad \textit{AB}, ut \textit{GC} ad \textit{GL}, seu ad \textit{DB}, seu}{%
\lemma{\textit{GL}}%
\Bfootnote{parallela ipsi \textit{AB}. %
\textit{(1)}~Ob %
\textit{(2)}~Jam (ob triangula \textit{ABC}, \textit{GLC} similia) %
\textit{(a)}~erit $AC : AB\; \squaredots\; GC : GL$ seu \textit{DB} seu \textit{CE} %
\textit{(b)}~est \textit{AC} %
ad \textit{AB} ut \textit{GC} ad \textit{GL} seu ad \textit{DB} seu~\textit{L}
\repeatlemma, \lbrack...\rbrack\ seu~\textit{E\textsuperscript{1}}}}
ad \textit{CE}. Rursus (ob triangula \textit{EGD}, \textit{ECH} similia)
\edtext{\textit{GC} ad \textit{CE},}{%
\lemma{\textit{GC} ad \textit{CE}}%
\Bfootnote{\textit{L}
\repeatlemma,~\textit{E\textsuperscript{1}}}}
ut \textit{DH} ad \textit{HE}. Ergo (a primis ad postrema) erit \textit{AC}
\edlabel{AE_1685_501-505_13a}%
\edtext{}{%
\xxref{AE_1685_501-505_13a}{AE_1685_501-505_13b}%
\lemma{ad \textit{AB}}%
\Bfootnote{%
ut \textit{DH} ad \textit{HE}. Quod ostendendum proponebatur. %
\bars{Possent ista extendi latius, si nunc immorari vacaret \textit{gestr.}} %
Patet autem ex %
\textit{(1)}~figura %
\textit{(2)}~demonstratione nihil~\textit{L}
\repeatlemma, \lbrack...\rbrack\ nihil~\textit{E\textsuperscript{1}}}}%
ad \textit{AB}, ut \textit{DH} ad \textit{HE}, quod ostendendum proponebatur.%
\edlabel{AE_1685_501-505_17}	%
\pend
\pstart
Patet autem ex demonstratione, nihil\edlabel{AE_1685_501-505_13b}
referre utrum angulus \textit{B} sit
\edtext{rectus,}{%
\lemma{rectus}%
\Bfootnote{;~\textit{L}
\repeatlemma,~\textit{E\textsuperscript{1}}}}
et proinde
\edtext{etiam \textit{DE}, basin trianguli \textit{ADE}, secare}{%
\lemma{etiam}%
\Bfootnote{%
\bars{rectam \textit{gestr.}} \textit{DE} basin trianguli \textit{ADE} secare~\textit{L}
\repeatlemma \lbrack...\rbrack\ secare~\textit{E\textsuperscript{1}}}}
basin \textit{BC} alterius trianguli \textit{ABC} (eandem laterum circa communem angulum summam habentis) ita ut segmenta \textit{CH}, 
\edtext{\textit{HB}, sint proportionalia ipsis \textit{DA}, \textit{AE}, lateribus trianguli}{%
\lemma{\textit{HB}}%
\Bfootnote{sint %
\textit{(1)}~in ratione %
\textit{(a)}~laterum trianguli \textit{DA}, \textit{AE}, %
\textit{(b)}~\textit{DA}, \textit{AE}, laterum %
\textit{(2)}~proportionalia ipsis %
\textit{(aa)}~lateribus \textit{DA}, \textit{AE}, %
\textit{(bb)}~\textit{DA}, \textit{AE}, lateribus trianguli~\textit{L}
\repeatlemma, \lbrack...\rbrack\ trianguli~\textit{E\textsuperscript{1}}}}
\textit{ADE} circa angulum
\edtext{triangulis}{%
\lemma{}%
\Bfootnote{triangulis \textit{fehlt~L}}}
communem \textit{A}. 
\pend
\pstart
\edtext{Potest etiam Lemma nostrum haberi pro accessione locorum Planorum. Locus enim punctorum, quibus bases omnium triangulorum, eandem summam laterum circa communem angulum habentium, in data ratione secantur, est recta.}{%
\lemma{}%
\Bfootnote{Potest \lbrack...\rbrack\ recta. \textit{fehlt~L}}}
\pend
\begin{edfigure}[1em][1em]
\includegraphics[scale=1]{%
gesamttex/edit_VIII-4/images/AE_1685_501-505_d2.pdf}
\edcaption{\lbrack\textit{Fig.~2; E\textsuperscript{1}}\rbrack}
\end{edfigure}
\pstart
\edlabel{AE_1685_501-505_3a}%
\edtext{}{%
\xxref{AE_1685_501-505_3a}{AE_1685_501-505_3b}%
\lemma{Superest}%
\Bfootnote{%
\textit{(1)}~ut %
\textit{(2)}~nunc ut \textls{veniamus ad solutionem casus de globo} \bars{\textit{I}~\textit{erg.}} %
\textls{duobus planis inclinatis} \bars{\textit{XC}, \textit{ZC} \textit{erg.}} %
\textls{ad se angulum} \lbrack...\rbrack\ \textls{incumbente ut determinemus quantum sit momentum quo} 
\textit{(a)}~unumquodque ex sin 
\textit{(b)}~\textls{unumquodque planum premitur.} %
\bars{Statim autem patet quod \textit{erg.}} %
ob \textls{duas causas} planum alterutrum ut \textit{XC} a globo~\textit{L}
\repeatlemma \lbrack...\rbrack\ \textls{planis}  \bars{\textit{XC}, \textit{ZC}, \textit{ändert Hrsg.}} \textls{ad horizontem} %
\bars{\textit{NO} \textit{ändert Hrsg.}} \textls{inclinatis,} \lbrack...\rbrack\ a globo~\textit{E\textsuperscript{1}}}}%
\edtext{Superest nunc, ut veniamus ad solutionem casus \textls{de globo}\indexS{globus duobus planis simul incumbens} \textit{I}
\edtext{(figur.~3)}{%
\lemma{(figur.~3)}%
\Cfootnote{Siehe hier das Diagramm \lbrack\textit{Fig.~2}\rbrack.}}
\textls{duobus planis} \lbrack\,\textls{\textit{XC}\,}\rbrack, %
\lbrack\,\textls{\textit{ZC}\,}\rbrack, %
\textls{ad horizontem} \lbrack\,\textls{\textit{NO}\,}\rbrack\ %
\textls{inclinatis, sed ad se invicem angulum rectum facientibus\indexS{plana inclinata angulum rectum facientia} simul incumbente, illi in \textit{F}, huic in \textit{H};} %
ut determinemus, quantum sit momentum, quo unumquodque planum premitur.\indexS{momentum quo planum inclinatum premitur}%
}{\lemma{}\Afootnote{%
\textit{Nur in L, am Rand von Bl.~2~r\textsuperscript{o}:} Hiehehr gehöhret die\textsuperscript{\lbrack a\rbrack} andere figur auß den \cite{01023}\textit{Actis} Novemb.~1684 \cite{01995}pag.~512%
\\[1.5ex]%
{\footnotesize \textsuperscript{\lbrack a\rbrack} die %
\textit{(1)}~figur auß %
\textit{(2)}~andere figur auß~\textit{L}}}}
\pend
\pstart
Statim autem patet, 
\edtext{(quod etiam ab Adm.\,R.\,P.~Kochanskio\indexN{\textls{Kochański}, Adam Adamandus 1631\textendash1700} in \cite{01023}\textit{Actis} Junii 1685 recte notatum video)}{%
\lemma{(quod \lbrack...\rbrack\ video)}%
\Cfootnote{\indexN{\textls{Kochański}, Adam Adamandus 1631\textendash1700}\textsc{Kochański}, \cite{01996}\textit{Consideratio speciminis}.}} 
globum\indexS{globus} in plano aliquo inclinato duplex exercere momentum,\indexS{momentum duplex} unum quo decliviter descendere 
\lbrack S.\;503\rbrack\ 
tendit,\indexS{momentum decliviter descendendi} alterum quo planum declive premit,\indexS{momentum planum declive premendi} 
quae duo simul absolutum seu totale gravis momentum\indexS{momentum gravis totale seu absolutum} constituunt. 
\edlabel{AE_1685_501-505_14a}%
Itaque in nostro casu \textls{ob duas causas} planum alterutrum ut \textit{XFC}, a globo\indexS{globus} %
\edlabel{AE_1685_501-505_3b}
\textit{I} premi intelligitur: \textls{prima} est quod globus\indexS{globus descendere tendens} \textit{I} descendere tendens in
\edtext{plani \textit{XFC}}{%
\lemma{}%
\Bfootnote{plani \textit{XFC}~\textit{fehlt~L}}}
linea \textit{FC},
\edtext{momento,}{%
\lemma{momento}%
\Bfootnote{\textit{L}
\repeatlemma,~\textit{E\textsuperscript{1}}}}
quod sit ad 
\edtext{totale,\indexS{momentum totale}}{%
\lemma{totale}%
\Bfootnote{\textit{L}
\repeatlemma,~\textit{E\textsuperscript{1}}}}
ut \textit{XN} ad
\edtext{\textit{XC} (quemadmodum demonstravimus) aget reliquo, (quod}{%
\lemma{\textit{XC}}%
\Bfootnote{(\textls{quemadmodum demonstravimus}) aget reliquo (quod~\textit{L}
\repeatlemma \lbrack...\rbrack\ (quod~\textit{E\textsuperscript{1}}}}
erit ad
\edtext{totale,}{%
\lemma{totale}%
\Bfootnote{\textit{L}
\repeatlemma,~\textit{E\textsuperscript{1}}}}
ut \textls{\textit{XC}\;\textendash\;}%
\edtext{\textls{\textit{XN},}}{%
\lemma{\textit{X\,N}}%
\Bfootnote{\textit{L}
\repeatlemma,~\textit{E\textsuperscript{1}}}}
\textls{ad \textit{XC}}) in ipsum planum declive
\edtext{\textit{XFC},}{%
\lemma{\textit{XFC}}%
\Bfootnote{\textit{L}
\repeatlemma,~\textit{E\textsuperscript{1}}}}
a quo
\edtext{sustentatur. Sed actio\indexS{actio} ista etsi vera et realis sit, est tamen alternativa et ob praesentiam alterius plani etiam sustentantis refracta, ut postea explicabitur. \textls{Secunda} causa premendi planum \textit{XFC} est,}{%
\lemma{sustentatur}%
\Bfootnote{%
ne perpendiculariter descendere possit. \textls{Secunda causa premendi} planum \textit{XC} est,~\textit{L}
\repeatlemma. Sed \lbrack...\rbrack\ est,~\textit{E\textsuperscript{1}}}}
quod globus\indexS{globus} \textit{I} descendens in altero plano \textit{ZHC}, momento quod sit ad
\edtext{totale,}{%
\lemma{totale}%
\Bfootnote{\textit{L}
\repeatlemma,~\textit{E\textsuperscript{1}}}} 
\textls{ut \textit{ZO}}
\edtext{\textls{ad \textit{ZC}}, premit}{%
\lemma{}%
\Bfootnote{%
\textls{ad \textit{Z\,C}} \bars{seu ad \textit{XC} \textit{gestr.}}, premit~\textit{L}}}
hoc ipso momento
\edtext{planum \textit{XFC} descensui suo declivi, in linea \textit{HC} affectato, se directe objiciens, atque ita hanc descensus vim excipiens. Porro}{%
\lemma{planum}%
\Bfootnote{%
\textit{(1)}~\textit{XC} %
\textit{(2)}~\textit{XFC} %
\bars{descensui suo declivi se directe objiciens atque \lbrack...\rbrack\ excipiens \textit{erg.}}. Porro~\textit{L}
\repeatlemma \lbrack...\rbrack\ Porro~\textit{E\textsuperscript{1}}}}
eadem ratiocinatio institui potest de
\edtext{plano altero \textit{ZHC}, etiam ex duabus illis premendi causis, et quidem \textls{ex prima causa} quae erat \textls{sustentatio} globi\indexS{globus} in ipso plano, in quo decliviter}{%
\lemma{plano}%
\Bfootnote{%
\textit{(1)}~\textit{ZC} altero %
\textit{(2)}~altero %
\textit{(a)}~\textit{ZC} %
\textit{(b)}~\textit{ZHC}, etiam ex duabus %
\bars{illis~\textit{erg.}} premendi causis, %
\bars{et quidem \textit{erg.}} \textls{ex prima} %
\textit{(aa)}~sustentationis propriae momento premitur a globo \textit{I}%
\textit{(bb)}~causa, quae erat sustentatio globi in ipso
\textit{(aaa)}~decurrere tendentis intelligi
\textit{(bbb)}~plano
\textit{(aaaa)}~\textit{ZC}
\textit{(bbbb)}~\textit{ZHC} decurrere tentantis
\textit{(cccc)}~dum in eo
\textit{(dddd)}~in quo decliviter~\textit{L}
\repeatlemma \lbrack...\rbrack\ decliviter~\textit{E\textsuperscript{1}}}}
descendere tendit,
\edtext{premitur hoc planum \textit{ZHC} momento\indexS{momentum}, quod sit ad totale momentum}{%
\lemma{premitur}%
\Bfootnote{%
\bars{hoc~\textit{erg.}} planum %
\bars{%
\textit{(1)}~\textit{ZC} %
\textit{(2)}~\textit{ZHC} \textit{erg.}} %
momento quod sit ad totale %
\textit{(1)}~, ut %
\textit{(2)}~momentum~\textit{L}
\repeatlemma \lbrack...\rbrack\ momentum~\textit{E\textsuperscript{1}}}}
\indexS{momentum globi}globi, \textls{ut \textit{ZC}\;\textendash\;}%
\edtext{\textls{\textit{ZO}, ad \textit{ZC}:} %
at vero \textls{ex altera} causa, quae erat \textls{exceptio} globi\indexS{globus}, %
dum scilicet planum \textit{ZHC}, ei in altero plano \textit{XFC}, decliviter descendere conanti, %
se directe objiciens, vim illam excipit, premitur planum \textit{ZHC} momento, quod}{%
\lemma{\textit{Z\,O}}%
\Bfootnote{%
\textls{ad} \textit{Z\,C}; %
\textls{ex altera} vero 
\textit{(1)}~exc %
\textit{(2)}~causa, \bars{quae erat \textit{erg.}} %
exceptio globi \bars{a plano \textit{ZC} \textit{erg.}} %
decliviter descendere conantis in altero plano \textit{XC}, %
premitur planum \textit{ZC} momento quod~\textit{L}
\repeatlemma \lbrack...\rbrack\ quod~\textit{E\textsuperscript{1}}}}
sit ad
\edtext{totale,}{%
\lemma{totale}%
\Bfootnote{\textit{L}
\repeatlemma,~\textit{E\textsuperscript{1}}}}
\textls{ut \textit{XN} ad \textit{XC}.}%
\edlabel{AE_1685_501-505_14b}%
\pend
\pstart
Et quidem causa 
\edtext{prima premendi planum \textit{XFC}, et secunda premendi planum \textit{ZHC}, quae}{%
\lemma{prima}%
\Bfootnote{%
\textit{(1)}~planum \textit{XC} %
\textit{(2)}~premendi planum \textit{XC}, et secunda premendi planum \textit{ZC} quae~\textit{L}
\repeatlemma \lbrack...\rbrack\ quae~\textit{E\textsuperscript{1}}}}
sunt dissimilares,
\edtext{simul}{\lemma{}\Bfootnote{simul~\textit{erg.~L}}}
constituunt totum momentum
\edtext{globi;\indexS{momentum globi totum}}{%
\lemma{globi}%
\Bfootnote{,~\textit{L}
\repeatlemma;~\textit{E\textsuperscript{1}}}}
globus enim premit
\edtext{planum \textit{XFC} quatenus ab eo sustentatur, et planum \textit{ZHC}}{%
\lemma{planum}%
\Bfootnote{%
\textit{XC} \lbrack...\rbrack\ planum \textit{ZC}~\textit{L}
\repeatlemma \lbrack...\rbrack\ \textit{ZHC}~\textit{E\textsuperscript{1}}}}
quatenus a priore non sustentatur. Similiter
\edtext{causa prima premendi planum \textit{ZHC}, et secunda premendi planum \textit{XFC}, licet dissimilares, constituunt totum momentum globi.\indexS{momentum globi totum} Causae}{%
\lemma{causa}%
\Bfootnote{%
\textit{(1)}~secunda premendi planum \textit{XC}, et %
\textit{(2)}~prima premendi planum \textit{ZC} et secunda premendi planum \textit{XC} licet dissimiles constituunt totum momentum globi. 
\href{\Scan{LH_35_11_14_002v}}%
{\lbrack2~v\textsuperscript{o}\rbrack} Causae~\textit{L}
\repeatlemma \lbrack...\rbrack\ Causae~\textit{E\textsuperscript{1}}}}
autem similes conjunctae, duae scilicet pressiones\indexS{pressio} ob sustentationem\indexS{sustentatio}, vel
\edlabel{AE_1685_501-505_9a}%
\edtext{}{%
\xxref{AE_1685_501-505_9a}{AE_1685_501-505_9b}%
\lemma{duae pressiones}%
\Bfootnote{%
\textit{(1)}~ad %
\textit{(2)}~ob exceptiones %
\bars{(quibus posterioribus compositis utitur 
\textit{(1)}~\textit{Specimen})
\textit{(2)}~doctissimus autor \textit{Speciminis} in \textit{Act.} Novemb. 1684) \textit{erg.}} %
momentum totale non integrant %
\bars{sumtae scilicet \lbrack...\rbrack\ partitionibus \textit{erg.}}. Verum cum %
\textit{(a)}~ita
\textit{(b)}~quatuor~\textit{L}
\repeatlemma \lbrack...\rbrack\ quatuor~\textit{E\textsuperscript{1}}}}%
duae pressiones\indexS{pressio} ob exceptionem,\indexS{exceptio} 
(quibus duabus posterioribus compositis utitur doctissimus 
\indexN{\textls{Vanni}, Giovanni Francesco 1638\textendash1709}Autor 
\cite{01995}\textit{Speciminis de momentis gravium} 
\edtext{in \cite{01023}\textit{Actis} Novembr.~1684}{\lemma{in \textit{Actis} Novembr. 1684}\Cfootnote{\textsc{Vanni},\indexN{\textls{Vanni}, Giovanni Francesco 1638\textendash1709} \cite{01995}\textit{Specimen}.}} 
exhibiti, qui casum hunc de duobus planis angulum rectum facientibus,\indexS{plana angulum rectum facientia} a se ingeniose excogitatum proposuit) %
momentum\indexS{momentum totale} totale non integrant, sumptae scilicet ex diversis momenti totalis partitionibus.\indexS{partitiones diversae momenti totalis} 
Verum cum quatuor%
\edlabel{AE_1685_501-505_9b}
premendi causis simul
\edtext{sumtis}{%
\lemma{sumtis}%
\Bfootnote{,~\textit{L}
\repeatlemma~\textit{E\textsuperscript{1}}}}
bis integretur momentum totale,\indexS{momentum totale}
\edtext{patet illas sic}{%
\lemma{patet}%
\Bfootnote{eas sic~\textit{L}
\repeatlemma illas sic~\textit{E\textsuperscript{1}}}}
absolute sumtas\indexS{causae absolute sumtae} non esse 
\edtext{compatibiles, nec \textls{cumulative},\indexS{causae cumulative sumtae} 
sed ut post dicam \lbrack S.\;504\rbrack\ tantum \textls{elective}\indexS{causae elective seu alternative compositae} sive alternative componendas:\indexS{causae alternative compositae}
alioqui effectus globi\indexS{globus} in plana major esset momento globi totalis absoluto.\indexS{momentum globi absolutum}\indexS{momentum globi totale seu absolutum}}{%
\lemma{compatibiles}%
\Bfootnote{%
\bars{nec componendas \textit{erg.}}, alioqui effectus globi %
\textit{(1)}~in %
\textit{(a)}~om %
\textit{(b)}~omnia %
\textit{(2)}~in plana, major esset momento globi \bars{totalis \textit{erg.}} absoluto.~\textit{L}
\repeatlemma, nec \lbrack...\rbrack\ absoluto.~\textit{E\textsuperscript{1}}}}
Cum vero manifestum 
\edtext{sit, duas}{%
\lemma{sit}%
\Bfootnote{%
\textit{(1)}~amb 
\textit{(2)}~duas~\textit{L}
\repeatlemma, duas~\textit{E\textsuperscript{1}}}}
semper causas in
\edlabel{AE_1685_501-505_5a}%
\edtext{}{%
\xxref{AE_1685_501-505_5a}{AE_1685_501-505_5b}%
\lemma{quolibet}%
\Bfootnote{%
\textit{(1)}~plano
\textit{(2)}~planorum %
in considerationem aequali
\textit{(a)}~jure %
\textit{(b)}~ratione venire \lbrack...\rbrack\ posse
adhibenda est \textls{regula alternativorum} \bars{seu \textls{concursuum} \textit{erg.}}, %
\textit{(aa)}~seu
\textit{(bb)}~hoc est utriusque %
\textit{(aaa)}~adhibendum %
\textit{(bbb)}~sumendum est dimidium
\textit{(ccc)}~momenti
\textit{(aaaa)}~sumendi
\textit{(bbbb)}~sumendum est \lbrack...\rbrack\ redit %
\bars{medium inter ipsa arithmeticum seu \textit{erg.}} dimidium~\textit{L}
\repeatlemma \lbrack...\rbrack\ dimidium~\textit{E\textsuperscript{1}}}}%
quolibet plano aequali ratione in considerationem venire debere, nec tamen integras retineri posse;
\edtext{adhibenda est \textls{regula alternativorum},\indexS{regula alternativorum} quae in jure accrescendi,\indexS{jus accrescendi} in aestimatione spei alea ludentium,\indexS{alea} aliisque casibus locum habet,}{%
\lemma{adhibenda \lbrack...\rbrack\ habet}%
\Cfootnote{%
In einigen Texten aus den Jahren 1686 und 1687 wendet Leibniz die \glqq\textls{regula alternativorum} (bzw. \glqq\textls{regula alternativorum} seu \textls{concursuum}\grqq\ in der Variante \textit{(b)} zu \refpassage{AE_1685_501-505_5a}{AE_1685_501-505_5b}) auf den schiefen Mehrkörperstoß an,
u.\,a. in der Aufzeichnung 
\cite{01398}\textit{De methodo alternorum, deque compositione motuum et virium in concursu obliquo} (\textit{LSB} VIII,~3 N.\;65\textsubscript{2}), unter ausdrücklicher Erwähnung des vorliegenden Aufsatzes (dort auf S.\;686),
und in dem Konzept \cite{02092}\textit{Motuum compositioni non est fidendum, praesertim in concursu obliquo} (\textit{LSB} VIII,~3 N.\;65\textsubscript{3}).}}
hoc est utriusque momenti\indexS{momentum} sumendum est dimidium, seu quod eodem redit, medium inter ipsa arithmeticum, sive dimidium%
\edlabel{AE_1685_501-505_5b}
summae ex ambobus. 
\pend
\pstart
Itaque si momentum
\edtext{globi}{%
\lemma{}%
\Bfootnote{%
globi~%
\textit{fehlt~L}}}
totale\indexS{momentum totale} sit unitas, cum duo momenta\indexS{momentum planum premendi} premendi
\edtext{planum \textit{XFC} sint \rule[-2ex]{0pt}{6.0ex}$\displaystyle\frac{XC-XN}{XC}$, et $\displaystyle\frac{ZO}{ZC}$ seu $\displaystyle\frac{ZO}{XC}$ (si ponantur \textit{XC} et \textit{ZC} aequales) sequitur momentum, quo}{%
\lemma{planum}%
\Bfootnote{%
\textit{XC} sint \rule[-2.2ex]{0pt}{6ex}$\displaystyle\frac{XC-XN}{XC}$ et $\displaystyle\frac{ZO}{ZC}$ sive $\displaystyle\frac{ZO}{XC}$ %
(quia \textit{ZC} et \textit{XC} aequales),
\textit{(1)}~itaque %
\textit{(2)}~sequitur momentum quo~\textit{L}
\repeatlemma \lbrack...\rbrack\ quo~\textit{E\textsuperscript{1}}}}
revera in hoc situ premitur
\edlabel{AE_1685_501-505_10a}%
\edtext{}{%
\xxref{AE_1685_501-505_10a}{AE_1685_501-505_10b}%
\lemma{planum}%
\Bfootnote{%
\textit{XC} erit \rule[-3ex]{0pt}{6.8ex}$\dfrac{XC - XN + ZO}{2XC}$, et similiter momentum quo
\textit{(1)}~premitur pl %
\textit{(2)}~revera premitur planum \textit{ZC}, erit \rule[-2.6ex]{0pt}{6ex}$\displaystyle\frac{XC - ZO + XN}{2XC}$ quorum %
\textit{(a)}~summa %
\textit{(b)}~duorum momentorum summa componit Momentum globi totale seu unitatem %
\bars{quod omnino \lbrack...\rbrack\ difficultas quae ex casu horum planorum nasci posse videbat \textit{erg.}}. %
In Numeris sit \textit{XN}, 3; \textit{ZO}, 4; \textit{XC}, 5,
\textit{(aa)}~planum \textit{XC} %
\textit{(bb)}~tunc planum \textit{XC} revera 
\textit{(aaa)}~$\displaystyle\frac{3}{5}$
\textit{(bbb)}~premitur 
\textit{(aaaa)}~trib
\textit{(bbbb)}~$\dfrac{3}{5}$
\textit{(cccc)}~a $\dfrac{3}{5}$ et planum \textit{ZC} a $\displaystyle\frac{2}{5}$,~\textit{L}
\repeatlemma \lbrack...\rbrack\ a $\dfrac{2}{5}$,~\textit{E\textsuperscript{1}}}}%
planum \textit{XFC}, fore dimid. \rule[-1.4ex]{0pt}{3.6ex}$\displaystyle\frac{XC - XN}{XC}\,+$ %
dimid. $\displaystyle\frac{ZO}{XC}$ seu $\displaystyle\frac{XC - XN\ \raisebox{0.05\baselineskip}{\rotatebox{90}{\dag}}\, ZO}{2XC}$, et similiter momentum\indexS{momentum quo planum premitur} quo revera premitur planum \textit{ZHC} erit dimid. 
\rule[-2.4ex]{0pt}{5ex}$\displaystyle\frac{ZC-ZO}{ZC}$ (seu dimid. $\displaystyle\frac{XC-ZO}{XC})\,+$ dimid. $\displaystyle\frac{XN}{XC}$ hoc est $\displaystyle\frac{XC-ZO\ \raisebox{0.35\baselineskip}{\rotatebox{-90}{\dag}}\,XN}{2XC}.$ 
quorum duorum momentorum revera duo plana prementium summa componit momentum globi totale\indexS{momentum globi totale} seu unitatem, quod omnino opus erat observari. 
\pend
\pstart
Atque adeo cessat difficultas, quae ex casu proposito nasci posse videbatur; quae tamen talis erat, ut omnino mereretur explicari: 
nec ideo minus obligati sumus \indexN{\textls{Vanni}, Giovanni Francesco 1638\textendash1709}Autori \cite{01995}\textit{Speciminis}, 
qui eam excogitavit, pariter atque \indexN{\textls{Kochański}, Adam Adamandus 1631\textendash1700}Adm.\,R.\,P.~Kochanskio, 
qui viam jam tum designavit, cui recte insistendo ad determinationem pressionis\indexS{pressio} %
cujusque plani perveniri poterat. 
In numeris \indexN{\textls{Kochański}, Adam Adamandus 1631\textendash1700}Kochanskianis
sit \textit{XN} vel \textit{CO}, 3. \textit{CN} vel \textit{ZO}, 4. \textit{XC} vel \textit{ZC}, 5. 
tunc planum \textit{XFC} revera premitur a $\displaystyle\frac{3}{5}$ et planum \textit{ZHC} a 
\rule[-2.6ex]{0pt}{6.0ex}$\displaystyle\frac{2}{5}$,%
\edlabel{AE_1685_501-505_10b}
momenti totalis.\indexS{momentum totale} Sin pro 3, 4, 5, numeri sint 5, 12, 13,
\edlabel{AE_1685_501-505_11a}%
\edtext{}{%
\xxref{AE_1685_501-505_11a}{AE_1685_501-505_11b}%
\lemma{tunc planum}%
\Bfootnote{%
\textit{XC} premitur \lbrack...\rbrack\ planum \textit{ZC} a 
\rule[-2ex]{0pt}{5.4ex}$\displaystyle\frac{3}{13}$, momenti totalis. %
\textit{(1)}~Quod fuerat %
\textit{(a)}~assignan %
\textit{(b)}~determinandum. Nec video qua alia %
\textit{(aa)}~ratione %
\textit{(bb)}~methodo omnium ratio habeatur. %
\textit{(aaa)}~Sic tamen ut in summa colligenda non impingatur
\textit{(bbb)}~Volo tamen
\textit{(ccc)}~Methodus autem ista per regulam alternativorum, %
etiam in aliis perplexis quaestionibus magni usus esse potest. %
\textit{(3)}~Ad examen~\textit{L}
\repeatlemma \lbrack...\rbrack\ examen~\textit{E\textsuperscript{1}}}}%
tunc planum \textit{XFC} premitur a 
\rule[-0ex]{0pt}{3.2ex}$\displaystyle\frac{10}{13}$, 
et planum \textit{ZHC} a $\displaystyle\frac{3}{13}$ momenti totalis.
\pend
\pstart
Ad examen%
\edlabel{AE_1685_501-505_11b}
autem regulae nostrae proderit etiam considerare, 
\edtext{et cum regula conferre solutiones casuum specialium aliunde notas: ut si}{%
\lemma{}%
\Bfootnote{%
et cum regula conferre solutiones casuum \bars{specialium \textit{erg.}} aliunde notas: ut si \textit{erg.~L}}}
\textit{XN} fiat aequalis nihilo,
\edtext{planum \textit{XFC} fiet horizontale,\indexS{planum horizontale} \textit{ZHC} verticale,\indexS{planum verticale} rectaeque \textit{XC} et \textit{ZO} aequales, ubi etiam ex regula recte prodibit}{%
\lemma{planum}%
\Bfootnote{%
\textit{XC}
\textit{(1)}~fieri
\textit{(2)}~fiet horizontale, \textit{ZC} verticale, et \textit{XC}, et \textit{ZO} aequales, ac proinde ex regula %
\bars{recte \textit{erg.}} %
\textit{(a)}~prodire
\textit{(b)}~prodibit~\textit{L}
\repeatlemma \lbrack...\rbrack\ prodibit~\textit{E\textsuperscript{1}}}}
totum momentum
\edtext{totale\indexS{momentum totale} pro premendo plano \textit{XFC}, et nihil pro premendo plano \textit{ZHC}. Si vero rectae \textit{XN}}{%
\lemma{totale}%
\Bfootnote{pro 
\textit{(1)}~\textit{XC}, nihil pro %
\textit{(2)}~premendo plano \textit{XC}, nihil pro premendo plano \textit{ZC}. Si vero \textit{XN}~\textit{L}
\repeatlemma \lbrack...\rbrack\ \textit{XN}~\textit{E\textsuperscript{1}}}}
et \textit{ZO} sumantur aequales,
\edtext{adeoque et anguli \textit{XCN}, \textit{ZCO} aequales, tunc per se manifestum est, utrumque}{%
\lemma{adeoque}%
\Bfootnote{%
\textit{(1)}~anguli \textit{XCN} %
\textit{(2)}~et anguli \textit{XCN}, \textit{ZCO} aequales, 
\textit{(a)}~manifestum est 
\bars{%
\textit{(1)}~aliunde %
\textit{(2)}~\bars{per \textit{streicht Hrsg.}}
\textit{erg.}}
\textit{(b)}~\bars{per se \textit{erg.}} manifestum est, utrumque~\textit{L}
\repeatlemma \lbrack...\rbrack\ utrumque~\textit{E\textsuperscript{1}}}}
plano eodem modo premi a globo, ac proinde dimidium
\edtext{momenti totalis unicuique impendi: quod}{%
\lemma{momenti}%
\Bfootnote{%
\textit{(1)}~unicuique %
\textit{(2)}~totalis unicuique impendi, quod~\textit{L}
\repeatlemma quod~\textit{E\textsuperscript{1}}}}
ipsum et ex regula nostra prodit. 
\pend
\pstart
Methodus\indexS{methodus per regulam alternativorum}
\edtext{autem qua}{%
\lemma{autem}%
\Bfootnote{%
,
\textit{(1)}~( %
\textit{(2)}~qua~\textit{L}
\repeatlemma qua~\textit{E\textsuperscript{1}}}}
usi sumus,
\edtext{per re\lbrack S.\;505\rbrack gulam alternativorum\indexS{regula alternativorum}, etiam}{%
\lemma{per}%
\Bfootnote{%
\textls{regulam alternativorum}, etiam~\textit{L}
\repeatlemma re\lbrack S.\;505\rbrack gulam alternativorum, etiam~\textit{E\textsuperscript{1}}}}
in aliis casibus perplexis enodandis magnum usum habere
\edlabel{AE_1685_501-505_12a}%
\edtext{}{%
\xxref{AE_1685_501-505_12a}{AE_1685_501-505_12b}%
\lemma{potest}%
\Bfootnote{%
.
\textit{(1)}~Et %
\textit{(2)}~Unum autem %
\textit{(a)}~add %
\textit{(b)}~adhuc admonendum est, (ne quis in captionem incidat %
\bars{nempe satis subtilem \textit{erg.}}), %
haec quidem
\textit{(aa)}~vera ess %
\textit{(bb)}~quae \lbrack...\rbrack\ determinavimus %
\textit{(aaa)}~vera es %
\textit{(bbb)}~tantum locum habere, cum %
\textit{(aaaa)}~ambo plana satis firma sunt ut resistere possint, %
\textit{(bbbb)}~unumquodque \lbrack...\rbrack\ possit %
\textit{(aaaaa)}~, alterutri conatuum alternativorum, %
\textit{(aaaaa-a)}~nam si %
\textit{(bbbbb-b)}~sed si maj %
\textit{(bbbbb)}~unicuique conatuum alternativorum, \bars{a quibus premitur, \textit{erg.}} separatim %
\textit{(aaaaa-a)}~spectat %
\textit{(bbbbb-b)}~sumto; sed si %
\textit{(aaaaa-aa)}~tam %
\textit{(bbbbb-bb)}~resistere possit %
\textit{(aaaaa-aaa)}~minori %
\textit{(bbbbb-bbb)}~majori ex ipsis %
\textit{(ccccc-cc)}~unum resistere quidem posset minori, vel \lbrack...\rbrack\ summae \bars{sed medio arithmetico \textit{erg.}}, non vero etiam majori %
\textit{(aaaaa-aaa)}~tunc cessante alternatione ipsa vi majoris conatus %
\textit{(bbbbb-bbb)}~alterum vero \lbrack...\rbrack\ esset ut %
\textit{(aaaaa-aaaa)}~ambobus %
\textit{(bbbbb-bbbb)}~unicuique~\textit{L}
\repeatlemma; quamvis \lbrack...\rbrack\ unicuique~\textit{E\textsuperscript{1}}}}%
potest; quamvis non sit pure Geometrica, sed quodammodo Metaphysica\indexS{regula simul Geometrica ac Metaphysica} simul; unde provisionalis censeri potest,
\edtext{donec idem demonstretur alia via, magis secundum vulgares notiones geometrica et rigorosa, quam ubi vacaverit, exhibere mihi reservo.}{%
\lemma{donec \lbrack...\rbrack\ reservo}%
\Cfootnote{Ein solcher Text ist nicht bekannt.}} 
\pend
\pstart
\edlabel{AE_1685_501-505_18a}%
Porro unum adhuc admonendum est, ne quis in captionem satis alioqui subtilem incidat, 
et ut appareat clarius, cur duos conatus\indexS{conatus alternativi} componendos vocaverimus alternativos: 
nempe haec quidem quae de pressione\indexS{pressio plani} cujusque plani determinavimus, 
tantum locum habent, cum unumquodque planum satis firmum est, ut resistere possit unicuique conatuum alternativorum\indexS{conatus alternativi} a quibus premitur, separatim sumto; 
sed si unum ex planis resistere quidem posset minori conatui\indexS{conatus urgens} se urgenti, 
vel etiam dimidiae amborum summae, non vero majori, alterum vero planum ita firmum esset, ut unicuique %
\edlabel{AE_1685_501-505_12b}
ex ambabus in se pressionibus posset resistere, 
\edtext{tunc eo ipso locum habebit absolute unus ex casibus alternativis \textls{prae reliquis}, ac proinde cessante regula alternativorum\indexS{regula alternativorum} aequaliter locum habentium, prius illud seu debilius planum vi ejus conatus,\indexS{conatus} cui impar est, superaretur.}{%
\lemma{tunc}%
\Bfootnote{%
\textit{(1)}~prius illud sed debilius %
\textit{(2)}~cessante alternativorum regula prius illud seu debilius %
\textit{(a)}~vi ejus conatus cui %
\textit{(aa)}~resiste %
\textit{(bb)}~impar est, sup
\textit{(b)}~planum vi ejus conatus cui impar est superaretur.~\textit{L}
\repeatlemma \lbrack...\rbrack\ superaretur.~\textit{E\textsuperscript{1}}}}
Quae quidem observatio singularis est, utilisque ut conatus\indexS{conatus alternativi} %
\edtext{naturae alternativi singuli, a conatu ob exitum}{%
\lemma{naturae}%
\Bfootnote{%
alternativus, a conatu denique post exitum~\textit{L}
\repeatlemma \lbrack...\rbrack\ exitum~\textit{E\textsuperscript{1}}}}
alterutro modo frustra
\edtext{tentatum,}{%
\lemma{tentatum}%
\Bfootnote{\textit{L}
\repeatlemma,~\textit{E\textsuperscript{1}}}}
ex ipsis %
\edtext{alternativis\indexS{conatus alternativi} sibi obstantibus resultante, distinguantur: nec ideo, exempli causa putetur planum \textit{XFC} premi absolute momento\indexS{momentum} 
\rule[-2.4ex]{0pt}{6ex}$\displaystyle\frac{ZO}{XC}$, quoniam}{%
\lemma{alternativis}%
\Bfootnote{%
\bars{sibi obstantibus \textit{erg.}} %
resultante distinguatur, nec ideo %
\bars{exempli causa \textit{erg.}} %
\textit{(1)}~petut %
\textit{(2)}~putatur %
\textit{(3)}~putetur %
\textit{(a)}~ambo plana pr %
\textit{(b)}~planum %
\textit{(aa)}~premi %
\textit{(bb)}~ut \textit{XFC} premi absolute momento globi \textit{I}, quod sit ad ejus gravitatem ut \textit{ZO} ad \textit{XC}, quoniam~\textit{L}
\repeatlemma \lbrack...\rbrack\ quoniam~\textit{E\textsuperscript{1}}}}
eo ipso momento\indexS{momentum} premitur %
\edlabel{AE_1685_501-505_4a}%
\edtext{}{%
\xxref{AE_1685_501-505_4a}{AE_1685_501-505_4b}%
\lemma{conditionaliter}%
\Bfootnote{%
, ita ut si ei resistere non posset frangeretur. %
Sunt autem haec quemadmodum paradoxa ita jucunda, et processum \lbrack...\rbrack\ mathematicum sed et quodammodo \protect{meta\greek{f}ysicum} illustrare apta.~\textit{L}
\repeatlemma; ita \lbrack...\rbrack\ apta.~\textit{E\textsuperscript{1}}}}%
conditionaliter; ita scilicet, ut frangendum sit, si ei resistere non possit. 
Nam in casu quo ei resistere potest, vi sua elastica\indexS{vis elastica} renitens prementi, 
onus in alterum planum rejicere tentat; quod ipsum si etiam satis firmum sit, ex conflictu ea demum aequissima distributio sequitur, quam exposuimus, ubi non ipsa quidem quatuor momenta,
\indexS{momentum} sed tamen ipsis proportionalia adhibentur, in quae momentum totale\indexS{momentum totale} distribuitur.%
\edlabel{AE_1685_501-505_18b}%
\pend
\pstart
Sunt autem haec, quemadmodum prima fronte paradoxa 
(ut scilicet planum aliquod dicatur certo aliquo momento\indexS{momentum alternative sive conditionaliter premendi} seu conatu\indexS{conatus alternative sive conditionaliter premendi} premi alternative sive conditionaliter)
ita consideratu jucunda, et processum naturae\indexS{processus naturae} non tantum mathematicum, sed et quodammodo Metaphysicum illustrare apta.%
\edlabel{AE_1685_501-505_4b}
\pend
